% ============================================================
% Section 14: Closing (Slides 100-101)
% ============================================================
\section{Closing}

% ----------------------------------------------------------
% Slide 100: Thank You / Questions
% ----------------------------------------------------------
\begin{frame}{Thank You --- Questions?}

  \begin{center}
    % Final mega-timeline (compact reprise)
    \visualplaceholder[5cm]{Main axis}
  \end{center}

  \vspace{0.8em}

  \begin{center}
    {\Large\bfseries\textcolor{digitalBlue}{What questions do you have?}}\\[1em]
    {\large\textcolor{gray}{Which mathematician's story surprised you the most?}}
  \end{center}

  \vspace{1em}

  \begin{center}
    \visualplaceholder[5cm]{Diagram}
  \end{center}

  \note{
    Open the floor for questions.
    Prompt: ``Which mathematician's story surprised you the most?''
    If no questions: highlight one or two stories that tend to resonate
    (Galois dying at 20, Ramanujan's self-teaching, Perelman's refusal,
    Mirzakhani's trailblazing).
    Thank the students for their attention.
  }
\end{frame}

% ----------------------------------------------------------
% Slide 101: Further Reading
% ----------------------------------------------------------
\begin{frame}{Further Reading}

  \begin{columns}[T]
    \begin{column}{0.55\textwidth}
      {\normalsize\bfseries Recommended books for curious minds:}

      \vspace{0.5em}
      \begin{enumerate}\setlength\itemsep{5pt}
        \item[\textcolor{digitalBlue}{\faBookOpen}]
          \textbf{Journey Through Genius}\\
          {\small William Dunham}\\
          {\scriptsize\itshape Great theorems told as human stories}

        \item[\textcolor{digitalBlue}{\faBookOpen}]
          \textbf{The Crest of the Peacock}\\
          {\small George Gheverghese Joseph}\\
          {\scriptsize\itshape Non-European roots of mathematics}

        \item[\textcolor{digitalBlue}{\faBookOpen}]
          \textbf{Mathematics in India}\\
          {\small Kim Plofker}\\
          {\scriptsize\itshape 500 BCE to 1800 CE --- the full Indian story}

        \item[\textcolor{digitalBlue}{\faBookOpen}]
          \textbf{A History of Mathematics}\\
          {\small Victor Katz}\\
          {\scriptsize\itshape The standard comprehensive textbook}

        \item[\textcolor{digitalBlue}{\faBookOpen}]
          \textbf{The Code Book}\\
          {\small Simon Singh}\\
          {\scriptsize\itshape From Caesar ciphers to quantum cryptography}
      \end{enumerate}
    \end{column}

    \begin{column}{0.42\textwidth}
      \visualplaceholder[5cm]{Book covers montage:\\
        Journey Through Genius,\\
        The Crest of the Peacock,\\
        Mathematics in India,\\
        A History of Mathematics,\\
        The Code Book}

      \vspace{0.8em}
      \begin{center}
        \visualplaceholder[3.5cm]{QR code linking to a curated\\
          online reading list}
      \end{center}

      \vspace{0.5em}
      \begin{center}
        {\small\bfseries\textcolor{digitalBlue}{Thank you!}}\\[0.3em]
        {\small Prof.\ J.\ Osterrieder $\cdot$ 2026}
      \end{center}
    \end{column}
  \end{columns}

  \note{
    Wrap-up slide. Mention each book briefly:
    \begin{itemize}
      \item Dunham: if you liked the stories, this book tells twelve great
            theorems as adventures.
      \item Joseph: if you want the full picture of non-European mathematics.
      \item Plofker: the definitive English-language account of Indian mathematics.
      \item Katz: the standard university textbook --- comprehensive and balanced.
      \item Singh: if you were interested in the cryptography section ---
            readable, exciting, covers everything from antiquity to quantum.
    \end{itemize}
    Final words: ``Mathematics is not over. The next great idea might come from you.
    Thank you.''
  }
\end{frame}
